\documentclass[11pt]{amsart}
\usepackage{amsmath,amssymb,amsthm}
\usepackage[margin=1.15in]{geometry}
\usepackage{booktabs}
\usepackage[hidelinks]{hyperref}
\hypersetup{pdftitle={One bit is optimal for cosmic coin flipping on at most nine bits},pdfauthor={Yitzchak Shmalo}}
\newtheorem{theorem}{Theorem}[section]
\newtheorem{lemma}[theorem]{Lemma}
\newtheorem{corollary}[theorem]{Corollary}
\theoremstyle{remark}
\newtheorem{remark}[theorem]{Remark}
\newcommand{\Z}{\mathbb{Z}}
\newcommand{\E}{\mathbb{E}}
\newcommand{\Prob}{\mathbb{P}}
\newcommand{\MAJ}{\mathrm{MAJ}}
\newcommand{\fl}{\operatorname{fl}}
\newcommand{\doi}[1]{\href{https://doi.org/#1}{doi:\nolinkurl{#1}}}
\newcommand{\arxiv}[1]{\href{https://arxiv.org/abs/#1}{arXiv:\nolinkurl{#1}}}
\setlength{\emergencystretch}{2.5em}

\title[Cosmic coin flipping on at most nine bits]{One bit is optimal for cosmic coin flipping on at most nine
bits}
\author{Yitzchak Shmalo}
\address{Einstein Institute of Mathematics, The Hebrew University of Jerusalem, Jerusalem, Israel}
\email{yspennstate@gmail.com}
\date{November 2026}
\subjclass[2020]{06E30, 68R05, 94A60, 65G50}
\keywords{non-interactive correlation distillation, cosmic coin flipping, noise stability, regular Boolean
functions, computer-assisted proof}

\raggedbottom
\begin{document}
\begin{abstract}
In the cosmic coin flipping problem of Mossel and O'Donnell, $k$ parties receive independent $\epsilon$-noisy copies of
$n$ uniform bits and each outputs a uniform bit, aiming to maximize the probability that all outputs agree. Mossel and
O'Donnell asked whether, for $2\le k\le 9$, the best antisymmetric protocol is for every party to output its first
bit. The cases $k\le5$ follow from work of Witsenhausen, Mossel and O'Donnell, and Yu. We prove that the answer is
positive for $6\le k\le 9$ whenever $n\le 9$. By symmetrization it suffices to compare the first-bit function with the
$319{,}124$ antisymmetric left-monotone functions on nine bits. After the substitution $t=2\epsilon/(1-2\epsilon)$ we
show that, for every other function, the coefficients of degrees $1$ to $nk-2$ of the numerator of the difference of
agreement probabilities are positive, using floating-point arithmetic with a rigorous error bound. For $n\le8$ an
exact computation confirms the inequalities.
\end{abstract}
\maketitle

\section{Introduction}

Let $x$ be uniform on $\{0,1\}^n$, and let $y^{(1)},\dots,y^{(k)}$ be obtained from $x$ independently by flipping each
coordinate with probability $\epsilon\in(0,1/2)$. Party $i$ sees only $y^{(i)}$ and outputs $f_i(y^{(i)})$, where
$f_i\colon\{0,1\}^n\to\{0,1\}$ is fixed in advance and balanced, that is, $|f_i^{-1}(1)|=2^{n-1}$. Mossel and
O'Donnell~\cite{MO05} asked for the protocol $(f_1,\dots,f_k)$ maximizing
\[
  P(f_1,\dots,f_k;\epsilon)=\Prob\bigl[f_1(y^{(1)})=\dots=f_k(y^{(k)})\bigr].
\]
A function $f$ is antisymmetric if $f(\bar x)=1-f(x)$, where $\bar x$ is the complement of $x$, and a protocol is
antisymmetric if all its functions are. Antisymmetric functions are balanced, and antisymmetry makes the common output
uniform. The same question on trees and with other noise models is the non-interactive correlation distillation
problem of Mossel, O'Donnell, Regev, Steif and Sudakov~\cite{MORSS06}; the case $k=2$ goes back to
Witsenhausen~\cite{W75}, and recent work of Kamath, Kumar and Manurangsi~\cite{KKM26} treats related two-party
conjectures. Li and M\'edard~\cite{LM21} study the same moments of noisy Boolean functions in the regimes of small
and large noise and of large moment order.

For $k=2$ and $k=3$ Mossel and O'Donnell proved that the best protocol has every party output a single bit of its
input~\cite[Theorem~1.3]{MO05}, and they showed that for fixed odd $n\ge3$ and fixed $\epsilon$ the majority of all
$n$ bits is optimal, and better than one bit, once $k$ is large enough~\cite[Theorem~1.5]{MO05}. Yu proved that among
all balanced functions the single bit maximizes the symmetric $\alpha$-stability for $2\le\alpha\le5$, in every
dimension and for every correlation~\cite[Corollary~2]{Yu23}; with $\alpha=k$ and correlation $1-2\epsilon$ this is the
case $k\le5$ of the question below. Write $\MAJ_r(x)$ for the majority of $x_1,\dots,x_r$ ($r$ odd); $\MAJ_1(x)=x_1$
is the first-bit function. For $x,y\in\{0,1\}^n$ write $x\le_L y$ if $\sum_{i\le m}x_i\le\sum_{i\le m}y_i$ for every
$m$, and call $f$ left-monotone if $f(x)\le f(y)$ whenever $x\le_L y$. Such functions are also called
regular Boolean functions. A signed permutation is a map $T(x)_j=x_{\pi(j)}\oplus b_j$ with $\pi$ a permutation of
the coordinates, $b\in\{0,1\}^n$ and $\oplus$ addition modulo $2$. The symmetrization theorem of Mossel and O'Donnell~\cite[Theorem~1.4]{MO05} says
that every protocol maximizing $P$, among all protocols or among antisymmetric ones, has $f_1=\dots=f_k=f$ with
$f\circ T$ left-monotone for a signed permutation $T$. When all parties use $f$,
\begin{equation}\label{eq:Pk}
  P_k(f;\epsilon):=P(f,\dots,f;\epsilon)=\E_x\bigl[p_f(x)^k+(1-p_f(x))^k\bigr],\qquad
  p_f(x)=\sum_{y}f(y)\,\epsilon^{d(x,y)}(1-\epsilon)^{n-d(x,y)},
\end{equation}
where $d$ is Hamming distance. For small $n$ the problem is therefore a comparison of finitely many polynomials in
$\epsilon$ on $(0,1/2)$. Mossel and O'Donnell counted $7$, $135$, $2470$ and $319{,}124$ antisymmetric left-monotone
functions for $n=5,7,8,9$, plotted the seven polynomials for $n=5$, and found that $\MAJ_1$ is best for $2\le k\le9$
while $\MAJ_3$ beats it on an interval of $\epsilon$ when $k=10$ or $11$~\cite[Fact~5.1, Proposition~5.2]{MO05}. They
posed the following.

\medskip\noindent\textbf{Open Problem 1} \cite{MO05}. Prove or disprove: for fixed $n$, $\epsilon$ and $2\le k\le 9$,
the best antisymmetric protocol is for all parties to use $\MAJ_1$.

\begin{theorem}\label{thm:main}
Let $1\le n\le 9$, $2\le k\le 9$ and $0<\epsilon<1/2$. For every antisymmetric left-monotone $f\colon\{0,1\}^n\to\{0,1\}$
other than $\MAJ_1$,
\[
  P_k(f;\epsilon)<P_k(\MAJ_1;\epsilon).
\]
\end{theorem}

\begin{corollary}\label{cor:op1}
Open Problem~1 has a positive answer for $n\le 9$: if $n\le 9$ and $2\le k\le 9$, every antisymmetric protocol that
maximizes $P(f_1,\dots,f_k;\epsilon)$ has all parties output the same function $x_i\oplus b$ for one coordinate $i$
and one $b\in\{0,1\}$.
\end{corollary}

\begin{proof}
A maximizer exists since there are finitely many protocols. By~\cite[Theorem~1.4]{MO05} it has the form
$f_1=\dots=f_k=f$ with $f\circ T$ left-monotone for some signed permutation $T$. A signed permutation commutes with
complementation of all coordinates and maps the uniform source and its noisy copies to a source and copies with the
same joint law, so $f\circ T$ is antisymmetric and $P_k(f\circ T;\epsilon)=P_k(f;\epsilon)$. Theorem~\ref{thm:main}
gives $f\circ T=\MAJ_1$. With $i=\pi^{-1}(1)$ we have $T(x)_i=x_1\oplus b_i$, so $f(z)=z_i\oplus b_i$ for every $z$. A function of fewer than $n$ bits is a
function of $n$ bits that ignores the rest, and appending ignored coordinates preserves antisymmetry,
left-monotonicity and the agreement probability, so the statement for $n=9$ contains those for $n<9$.
\end{proof}

For $k=2,3$ the theorem is contained in~\cite[Theorem~1.3]{MO05}, and for $k=4,5$ the optimality of $\MAJ_1$ follows
from~\cite[Corollary~2]{Yu23}, which covers all balanced functions in every dimension; the new cases are
$6\le k\le 9$. The computation below covers $k\le5$ as well and gives the strict inequality there.
For $n\ge5$ the range of $k$ cannot be extended, since for $n=5$ and $k=10$ the function $\MAJ_3$ beats $\MAJ_1$ at
$\epsilon=0.26$~\cite[Proposition~5.2]{MO05}. The proof of Theorem~\ref{thm:main} is a computation, organized so that
every step is either exact or carries a proved error bound. Section~\ref{sec:poly} reduces each comparison to the
positivity of the coefficients of one polynomial, Section~\ref{sec:float} explains why floating-point coefficients
suffice, and Section~\ref{sec:comp} describes the enumeration and the results, together with an exact check for
$n\le8$.

\section{Agreement probabilities as polynomials}\label{sec:poly}

Put $t=2\epsilon/(1-2\epsilon)$, an increasing bijection from $(0,1/2)$ onto $(0,\infty)$. Then
\[
  \epsilon=\frac{t}{2(1+t)},\qquad 1-\epsilon=\frac{2+t}{2(1+t)},\qquad
  \epsilon^{d}(1-\epsilon)^{n-d}=\frac{t^{d}(2+t)^{n-d}}{(2(1+t))^{n}} .
\]
For $f\colon\{0,1\}^n\to\{0,1\}$ and $x\in\{0,1\}^n$ let
\[
  q_{f,x}(t)=\sum_{y}f(y)\,t^{d(x,y)}(2+t)^{n-d(x,y)},\qquad R_{f,k}(t)=\sum_{x}q_{f,x}(t)^k .
\]
Each $q_{f,x}$ has degree at most $n$ and nonnegative integer coefficients, and $p_f(x)=q_{f,x}(t)/(2(1+t))^n$. If $f$
is antisymmetric then $p_f(\bar x)=1-p_f(x)$, because $d(\bar x,y)=d(x,\bar y)$ and $f(\bar y)=1-f(y)$, where $\bar y$ is the complement of $y$. So the two
sums in~\eqref{eq:Pk} are equal, and
\begin{equation}\label{eq:PR}
  P_k(f;\epsilon)=\frac{2^{1-n}\,R_{f,k}(t)}{(2(1+t))^{nk}} .
\end{equation}

\begin{lemma}\label{lem:ends}
For balanced $f$ the coefficients of $t^0$, $t^{nk-1}$ and $t^{nk}$ in $R_{f,k}$ do not depend on $f$.
\end{lemma}

\begin{proof}
Only $y=x$ contributes to the constant term of $q_{f,x}$, so $q_{f,x}(0)=2^nf(x)$ and the constant term of $R_{f,k}$
is $2^{nk}|f^{-1}(1)|=2^{nk+n-1}$. Each $t^{d}(2+t)^{n-d}$ is monic of degree $n$ with coefficient $2(n-d)$ at
$t^{n-1}$. Hence the leading coefficient of every $q_{f,x}$ is $|f^{-1}(1)|=2^{n-1}$, and its coefficient at $t^{n-1}$
is $\sum_y f(y)\,2(n-d(x,y))$. It follows that the coefficient of $t^{nk}$ in $R_{f,k}$ is $2^n2^{(n-1)k}$ and that
of $t^{nk-1}$ is
\[
  k\,2^{(n-1)(k-1)}\sum_{y}f(y)\sum_{x}2\bigl(n-d(x,y)\bigr)=k\,2^{(n-1)(k-1)}\cdot2^{n-1}\cdot n\,2^{n},
\]
since $\sum_x d(x,y)=n2^{n-1}$ for every $y$.
\end{proof}

Let $D_{f,k}=R_{\MAJ_1,k}-R_{f,k}$. By Lemma~\ref{lem:ends}, $D_{f,k}$ has zero constant term and degree at most
$nk-2$. By~\eqref{eq:PR}, $P_k(\MAJ_1;\epsilon)-P_k(f;\epsilon)$ has the sign of $D_{f,k}(t)$. Hence:

\begin{corollary}\label{cor:coef}
Let $f$ be antisymmetric and $nk\ge3$. If the coefficients of $t^1,\dots,t^{nk-2}$ in $D_{f,k}$ are all
positive, then
$P_k(f;\epsilon)<P_k(\MAJ_1;\epsilon)$ for every $\epsilon\in(0,1/2)$.
\end{corollary}

The condition is sufficient, not necessary. It is the case of Descartes' rule of signs with no sign change, and it
fails, for example, for $f=\MAJ_3$ with $n=5$ and $k=10$, where the difference changes sign. For the values in
Theorem~\ref{thm:main} it always holds.

\section{Floating-point coefficients with a proved error bound}\label{sec:float}

The coefficients of $R_{f,k}$ are integers of up to about $50$ decimal digits when $n=k=9$, and computing them
exactly for all $319{,}124$ functions is slow. Since every quantity entering $R_{f,k}$ is nonnegative, ordinary
double-precision arithmetic has a relative error that can be bounded in advance, and a coefficient of $D_{f,k}$ is
certified positive once it exceeds that bound.

We use the standard model of floating-point arithmetic: for $a,b$ representable and $\circ\in\{+,\times\}$, the computed
result is $\fl(a\circ b)=(a\circ b)(1+\delta)$ with $|\delta|\le u=2^{-53}$, as long as neither overflow
nor underflow occurs~\cite[Section~2.2]{Hig02}. Since $q_{f,x}(1)\le\sum_y3^{n-d(x,y)}=4^n$, every coefficient of $R_{f,k}$ is at
most $2^n4^{nk}<10^{52}$, far below the overflow threshold, and every nonzero value computed below, the thresholds included, exceeds
$10^{-15}$, far above the underflow threshold. If $|\delta_i|\le u$ and $mu<1$, then
$\prod_{i=1}^m(1+\delta_i)=1+\theta$ with $|\theta|\le\gamma_m=mu/(1-mu)$~\cite[Lemma~3.1]{Hig02}.

\begin{lemma}\label{lem:nonneg}
Assume the model above and $mu<1$. Let a number $A$ be computed from exactly representable nonnegative inputs by additions and multiplications only. Each
occurrence of a term in the fully unfolded expansion of $A$ as a sum of products of inputs is then multiplied by its
own product of factors $1+\delta_i$, one for each rounding that affects it, counted with multiplicity. If no occurrence
carries more than $m$ such factors, the computed value is $\hat A=A(1+\theta)$ with $|\theta|\le\gamma_m$.
\end{lemma}

\begin{proof}
Each term $T_j\ge0$ of $A$ appears in $\hat A$ as $T_j(1+\theta_j)$ with $|\theta_j|\le\gamma_m$ by the bound above,
so $\hat A-A=\sum_jT_j\theta_j$ and $|\hat A-A|\le\gamma_m\sum_jT_j=\gamma_mA$.
\end{proof}

The arithmetic is IEEE binary64 with rounding to nearest. We compute the coefficient vector of $q_{f,x}$ from the
expansion of $t^d(2+t)^{n-d}$ as a matrix product. Every input, product and partial sum in it is a nonnegative integer
at most $q_{f,x}(1)\le4^n\le2^{18}$, so it is computed exactly, in any order of summation and with or without fused
multiply-add. We then form $q_{f,x}^k$ by $k-1$ convolutions with the exact vector $q_{f,x}$. Each coefficient of a
convolution is a sum of at most $n+1$ products, accumulated one at a time, so an occurrence of a term acquires at most
$n+2$ rounding factors per convolution. Summing over the $2^n$ points $x$ adds at most $2^n-1$ more, in any order.
Hence every coefficient $A$ of $R_{\MAJ_1,k}$ and $B$ of $R_{f,k}$ is computed with relative error at most $\gamma_m$,
where $m=(k-1)(n+2)+2^n+1$. If $\hat A=A(1+\theta_A)$ and $\hat B=B(1+\theta_B)$, then
\[
  A-B\ \ge\ \frac{\hat A}{1+\gamma_m}-\frac{\hat B}{1-\gamma_m},
\]
which is positive as soon as $\hat A-\hat B>\gamma_m(\hat A+\hat B)$. Our program computes $\gamma=4mu/(1-mu)$ and
accepts a coefficient only if $\fl(\hat A-\hat B)>\fl(2\gamma\,\fl(\hat A+\hat B))$. The computed $\gamma$ is at least
$4\gamma_m(1-u)/(1+u)$, so the computed threshold is at least $8\gamma_m(\hat A+\hat B)(1-u)^3/(1+u)$, and acceptance
gives $\hat A-\hat B>8\gamma_m(\hat A+\hat B)(1-u)^3/(1+u)^2>\gamma_m(\hat A+\hat B)$. For $n=k=9$ we have $m=601$ and
$8\gamma_m<5.4\cdot10^{-13}$.

\section{The computation}\label{sec:comp}

This section describes how the functions are listed, gives the outcome of the certificate, and reports two checks
made by different methods.

\subsection*{Enumeration}
Write $\uparrow\!x=\{y:x\le_L y\}$ and $\downarrow\!x=\{y:y\le_L x\}$. Complementation reverses $\le_L$, since the
prefix sums of $\bar x$ are $m-\sum_{i\le m}x_i$; so it maps $\uparrow\!x$ onto $\downarrow\!\bar x$. The antisymmetric
left-monotone functions are listed by a depth-first search. A state is a pair of disjoint sets $(I,O)$, with $I$
up-closed and $O=\bar I:=\{\bar x:x\in I\}$. At each step the least undecided point $x$ is set either to $1$, adding $\uparrow\!x$ to $I$
and $\downarrow\!\bar x$ to $O$, or to $0$, adding $\uparrow\!\bar x$ to $I$ and $\downarrow\!x$ to $O$. A branch
whose two sets meet is abandoned. Every antisymmetric left-monotone function is reached exactly once, at the leaf
whose choices agree with it. The search returns $1,1,2,3,7,21,135,2470,319124$ functions for $n=1,\dots,9$. This
agrees with the counts of Mossel and O'Donnell~\cite{MO05} and with the number of self-dual regular Boolean functions
computed by Knuth (sequence A109456 of~\cite{OEIS}; see also Hausmann~\cite{Hau15}). As a further check, for $n=7$ we
filtered the $1{,}422{,}564$ self-dual monotone functions for left-monotonicity by a separate program and obtained the
same $135$.

\subsection*{Results}
For $n\le2$ the only antisymmetric left-monotone function is $\MAJ_1$. For $3\le n\le 9$, every $k$ with
$2\le k\le 9$ and every antisymmetric left-monotone $f\ne\MAJ_1$, all coefficients
of $t^1,\dots,t^{nk-2}$ in $D_{f,k}$ passed the test of Section~\ref{sec:float}. By Corollary~\ref{cor:coef} this
proves Theorem~\ref{thm:main}. Table~\ref{tab:margin} shows how far the test is from failing. It lists the smallest
value of $(\hat A-\hat B)/(\hat A+\hat B)$ over all functions, all $k$ and all middle coefficients; the relative
threshold $2\gamma$ is below $6\cdot10^{-13}$.

\begin{table}[tbp]
\centering
\small
\setlength{\tabcolsep}{4pt}
\begin{tabular}{lccccccc}
\toprule
$n$ & 3 & 4 & 5 & 6 & 7 & 8 & 9\\\midrule
functions & 2 & 3 & 7 & 21 & 135 & 2470 & 319124\\
margin & $7.9\cdot10^{-3}$ & $4.3\cdot10^{-3}$ & $1.9\cdot10^{-3}$ & $7.6\cdot10^{-4}$ &
$3.0\cdot10^{-4}$ & $1.2\cdot10^{-4}$ & $4.9\cdot10^{-5}$\\\bottomrule
\end{tabular}
\medskip
\caption{The margin is the smallest value of $(\hat A-\hat B)/(\hat A+\hat B)$, where $\hat A-\hat B$ is the computed
coefficient of $t^j$ in $D_{f,k}$, over $1\le j\le nk-2$, all antisymmetric left-monotone $f\ne\MAJ_1$ and
$2\le k\le 9$. In every case it is attained
at $k=2$. For $n\le2$ the function $\MAJ_1$ is the only one.}
\label{tab:margin}
\end{table}

As a check on the method, we applied the same program to $n=5$ and $k=10,11,12$. It certifies $4$, $4$ and $2$ of the
six functions other than $\MAJ_1$. At $k=10$ one of the two rejected functions is $\MAJ_3$, which beats $\MAJ_1$ on an
interval of $\epsilon$, in agreement with~\cite[Fact~5.1]{MO05}. The other is $\MAJ_5$, which lies below $\MAJ_1$ on
all of $(0,1/2)$ although its difference with $\MAJ_1$ has a negative coefficient. At $k=11$ the rejected functions
are again $\MAJ_3$ and $\MAJ_5$, and both beat $\MAJ_1$ somewhere.

\subsection*{\texorpdfstring{Exact check for $n\le8$}{Exact check for n <= 8}}
For $3\le n\le 8$ and $4\le k\le 9$ we also computed $P_k(\MAJ_1;\epsilon)-P_k(f;\epsilon)$ in exact rational
arithmetic, as a polynomial in $\epsilon$, with SymPy~\cite{SymPy}. We wrote it as $\epsilon^a(1-2\epsilon)^bg(\epsilon)$,
with $a,b\ge0$ integers and $g$ a polynomial not vanishing at $0$ or $1/2$, substituted $\epsilon=t/(2(1+t))$ in $g$ and cleared the denominator. For every $f\ne\MAJ_1$ the coefficients of the resulting
polynomial in $t$ were nonnegative and not all zero, which proves $P_k(f;\epsilon)<P_k(\MAJ_1;\epsilon)$ on $(0,1/2)$
independently of Section~\ref{sec:float}. Separately, we evaluated
$(P_k(\MAJ_1;\epsilon)-P_k(f;\epsilon))/(\epsilon(1-2\epsilon)^2)$ in double precision at $4000$ equally spaced points
of $(0,1/2)$ for $n=7,8$ and $2\le k\le 9$; for $n=7$ the functions were obtained by a second program, which filters
the self-dual monotone functions. The smallest value found was $1.6\cdot10^{-2}$ for $n=7$ and $8.2\cdot10^{-3}$ for
$n=8$, both at $k=9$ near $\epsilon=1/2$.

\subsection*{Larger values of the number of bits}
For $n=10$ there are $1{,}214{,}554{,}343$ self-dual regular functions~\cite{OEIS}, and the same computation would take
several thousand times longer. From $n=5$ on, the relative margins in Table~\ref{tab:margin} shrink by a factor of
about $2.5$ for each added bit. They give no evidence against a positive answer to Open Problem~1 for every $n$, but they do not
suggest a proof either.

\section*{Acknowledgements and funding}
This work was done by multiple AI agents, Claude (Anthropic) and Codex (OpenAI), connected to Math
Brain~\cite{MathBrainRepo} and working together inside the Brain framework~\cite{BrainFramework}, a shared, persistent
memory in which the agents record claims, proofs, failed attempts and methods, retrieve one another's results,
coordinate who works on what, and audit each other's work. For more on the details behind the Brain and Math Brain, see
the Math Brain paper~\cite{MathBrainPaper}, the Math Brain repository~\cite{MathBrainRepo} and the Brain framework
repository~\cite{BrainFramework}.

The research presented in this paper was supported by the European Research Council (ERC) under the European Union's
Horizon Europe research and innovation programme (grant agreement No.~101041711), by the Simons Foundation as part of
the Collaboration on the Mathematical and Scientific Foundations of Deep Learning, by Heights Labs, by Convex Nexus
Capital, by the Israel Science Foundation (grant number 2258/19), and by the Israel Science Foundation (ISF Grant
4101/25). Computational resources were provided in part by Google Cloud through its education credits program (credit
530550103).

\section*{Data availability}
The programs that enumerate the functions and certify the coefficients, the exact and numerical checks, and their
output are available with the \LaTeX{} source of the paper at\\
\url{https://github.com/yspennstate/cosmic-coin-flipping-nine-bits/tree/v1.0}.
The computations were run with Python 3.13.2, NumPy 2.2.3 (Accelerate BLAS) and SymPy 1.13.1.

The paper is part of the Math Brain project, whose public repository is given in~\cite{MathBrainRepo}.

\section*{Declaration of competing interest}
The author holds equity in Heights Labs and in Convex Nexus Capital and has a management role at Convex Nexus Capital.
These entities are included in the funding declaration above.

\section*{Declaration of generative AI and AI-assisted technologies}
Claude agents (Anthropic) designed and ran the computation, proved Lemmas~\ref{lem:ends} and~\ref{lem:nonneg}, and
prepared the manuscript; Codex agents (OpenAI) refereed the proofs and the programs. The problem, the reduction to one
left-monotone function, the cases $k=2,3$, the function counts and the analysis of $n=5$ are due to Elchanan Mossel and
Ryan O'Donnell~\cite{MO05}, the case $k=2$ goes back to Hans Witsenhausen~\cite{W75}, and the cases $k=4,5$ follow from
Lei Yu's theorem on symmetric stability~\cite{Yu23}; the rounding-error bounds are Nicholas Higham's~\cite{Hig02}.
Corollary~\ref{cor:coef} is an elementary sufficient condition. New here are the certified computation and
Theorem~\ref{thm:main} for $6\le k\le9$. The author directed the work and is responsible for the content of the
paper.

\begin{thebibliography}{99}
\bibitem{Hau15} J.-C. Hausmann, \emph{Counting polygon spaces, Boolean functions and majority games}, 2015,
\arxiv{1501.07553}.
\bibitem{Hig02} N.~J. Higham, \emph{Accuracy and Stability of Numerical Algorithms}, 2nd ed., SIAM, Philadelphia,
2002, \doi{10.1137/1.9780898718027}.
\bibitem{KKM26} P.~Kamath, R.~Kumar and P.~Manurangsi, \emph{When majority fails: tight bounds for correlation
distillation conjectures}, 2026, \arxiv{2604.06590}.
\bibitem{LM21} J.~Li and M.~M\'edard, \emph{Boolean functions: noise stability, non-interactive correlation
distillation, and mutual information}, IEEE Trans. Inform. Theory \textbf{67} (2021), 778--789,
\doi{10.1109/TIT.2020.3041028}, \arxiv{1801.04462}.
\bibitem{MO05} E.~Mossel and R.~O'Donnell, \emph{Coin flipping from a cosmic source: on error correction of truly
random bits}, Random Structures Algorithms \textbf{26} (2005), 418--436, \doi{10.1002/rsa.20062}, \arxiv{math/0406504}.
\bibitem{MORSS06} E.~Mossel, R.~O'Donnell, O.~Regev, J.~E. Steif and B.~Sudakov, \emph{Non-interactive correlation
distillation, inhomogeneous Markov chains, and the reverse Bonami--Beckner inequality}, Israel J. Math. \textbf{154}
(2006), 299--336, \doi{10.1007/BF02773611}, \arxiv{math/0410560}.
\bibitem{OEIS} OEIS Foundation, \emph{The On-Line Encyclopedia of Integer Sequences}, sequence A109456, number of
Boolean functions of $n$ variables that are self-dual and regular, contributed by D.~E. Knuth, accessed November 2026,
\url{https://oeis.org/A109456}.
\bibitem{SymPy} SymPy Development Team, \emph{SymPy: a Python library for symbolic mathematics}, version 1.13.1, 2024,
\url{https://www.sympy.org}.
\bibitem{W75} H.~S. Witsenhausen, \emph{On sequences of pairs of dependent random variables}, SIAM J. Appl. Math.
\textbf{28} (1975), 100--113, \doi{10.1137/0128010}.
\bibitem{Yu23} L.~Yu, \emph{On the $\Phi$-stability and related conjectures}, Probab. Theory Related Fields
\textbf{186} (2023), 1045--1080, \doi{10.1007/s00440-023-01209-5}, \arxiv{2104.08740}.
\bibitem{MathBrainRepo} Y.~Shmalo, \emph{Math Brain}, public software repository, snapshot of 14 November 2026,
\url{https://github.com/yspennstate/math-brain-public/tree/8c9a6bd478c4ec4082347151b7ae16db50c917dd}.
\bibitem{MathBrainPaper} Y.~Shmalo, \emph{Math Brain: shared, auditable memory for multi-agent mathematical research},
manuscript dated 26 October 2026, Zenodo release 12 November 2026,
\url{https://github.com/yspennstate/math-brain-public/tree/main/papers/math-brain-2026-10-24},
\doi{10.5281/zenodo.24150837}.
\bibitem{BrainFramework} Y.~Shmalo, \emph{Brain Framework}, public software repository, snapshot of 13 November 2026,
\url{https://github.com/yspennstate/brain-framework-public/tree/5629e644cdd510173ae97e9aa92642dabd22dc57}.
\end{thebibliography}

\end{document}
